\subsection{Rule Consolidation}
\label{sec:consolidate}

Rule generation typically produces a large set of overlapping and overly specific rules.
Consolidation is therefore required to select a compact subset that generalizes well.

%Humans do not retain every exception they encounter; instead, they consolidate experience by pruning, merging, and abstracting rules to maintain a compact, reusable knowledge base. We formalize this intuition using the Minimum Description Length (MDL) principle \citep{rissanen1978mdl,grunwald2007mdl}, which favors rule libraries that balance explanatory power with conciseness by compressing both the rules themselves and the remaining failures they explain. 
%Effective use of such rules also requires accurate retrieval at inference time: relevant rules must be identified reliably based on the current query and available tools, while ignoring incidental linguistic variation in how rules are phrased. We therefore represent rules in a dual form—a natural-language version for interpretability and a symbolic representation that abstracts away surface-level wording, enabling principled consolidation and efficient, structure-aware retrieval.
%Ablation studies show that both components of our design are essential: MDL-guided consolidation yields more compact and effective rule libraries than prompting-based consolidation, and symbolic, structure-aware retrieval consistently outperforms retrieval based on semantic similarity over raw natural-language rules.

We consider two approaches: a prompt-based heuristic that merges rules using an LLM, and a data-driven alternative based on the Minimum Description Length (MDL) principle.


\subsubsection{Prompt-Based}

A simple consolidation strategy is to prompt an LLM to merge or rewrite candidate rules,
as in AutoRule~\citep{wang2025autorule}. While this can reduce redundancy at the textual
level, it is heuristic and does not account for how rules affect agent behavior on actual
failure cases. This motivates a consolidation objective that explicitly trades off rule
complexity against empirical effectiveness.


\subsubsection{MDL-Guided}
\label{sec:mdl}

%Unlike neural models with fixed capacity, our adaptation takes the form of a symbolic rule library whose size and content are explicitly optimized. 
We adopt the \emph{Minimum
Description Length} (MDL) principle, which favors models that compress observed data while
remaining concise~\citep{rissanen1978modeling,Grunwald07MDL}. In our setting, MDL provides
a principled objective for selecting a compact rule set that corrects as many failures as
possible without unnecessary complexity.

\paragraph{Objective.}
Let \( H \subseteq \mathcal{R}^{\text{sym}} \) denote a candidate rule library. We minimize
\[
\mathrm{MDL}(H) = L(H) + L(D \mid H),
\]
where \( L(H) \) penalizes rule complexity and \( L(D \mid H) \) measures how well the rules
correct observed failures.

\paragraph{Model Cost.}
We define a length-based prior over rule libraries,
\[
P(H) \propto \exp\!\Bigl(-\alpha \sum_{r \in H} \ell(r)\Bigr),
\]
where \( \ell(r) \) denotes the \textbf{symbolic token length}
\footnote{Symbolic token length provides a canonical measure of rule complexity by
abstracting away incidental natural-language phrasing.}
of rule \( r \) and \( \alpha > 0 \)
controls the strength of regularization. This corresponds to a maximum-entropy prior under
a constraint on expected total rule length~\citep{jaynes1957maxent}. Taking the negative
log-likelihood yields the model code length
\[
L(H) = -\log P(H) = \alpha \sum_{r \in H} \ell(r) + \log Z(\alpha),
\]
where \( Z(\alpha) \) is constant for a fixed candidate pool and is therefore omitted during
optimization.

\paragraph{Data Cost.}
Given failure cases \( D = \{(x_i, \mathcal{S}_i, \tau_i^-, \tau_i^*)\}_{i=1}^n \), we define
\( a_i(H) = 1 \) if injecting \( H \) corrects failure \( i \). Let
\( k_H = \sum_i a_i(H) \). We model correction outcomes with a Bernoulli likelihood and
use the refined MDL plug-in code
\[
L(D \mid H) = -\bigl[k_H \log \hat{p}_H + (n-k_H)\log(1-\hat{p}_H)\bigr],
\]
where \( \hat{p}_H = k_H/n \). This term directly rewards rule sets that correct more failures.

\paragraph{Greedy Consolidation.}
We minimize \( \mathrm{MDL}(H) \) starting from the full rule pool and applying local edits
that strictly reduce the objective. We consider two operations: \emph{pruning} a rule, and
\emph{generalizing} a tool-specific rule to a category-level rule. Edits are accepted only
when the reduction in model cost outweighs any increase in data cost. This greedy procedure
terminates when no local edit improves the MDL objective.

In practice, consolidation scales linearly with the size of the candidate rule pool, as each iteration evaluates a constant set of local edits per rule. While this greedy procedure does not guarantee a globally optimal solution and may miss beneficial non-local edits, it follows standard practice in rule learning and MDL-based pruning, where local edits are commonly used to remove redundancy and over-specific rules and are empirically effective in improving generalization~\citep{Quinlan87Simplifying,Furnkranz97Pruning,Grunwald07MDL}. Exploring more global or approximate inference strategies for rule consolidation is an interesting direction for future work.


\paragraph{MDL Behavior.}
Although the term \( L(D \mid H) \) could be minimized by correcting no failures,
this degenerate case is never reached in practice. Consolidation proceeds greedily from a
high-performing initial rule set; in this regime, removing a rule that reduces the number
of corrected failures sharply increases \( L(D \mid H) \), outweighing the modest reduction
in model cost. Such edits are therefore rejected.

\paragraph{Choice of $\alpha$.}
We select the regularization strength \( \alpha \) by choosing the value that yields the best
end-to-end accuracy after consolidation on the failure set. Since each rule is evaluated on
failures beyond its originating instance, this selection already reflects cross-sample
generalization. Performance is stable across a wide range of \( \alpha \), indicating that
MDL acts as a robust regularizer.
